\section{Tipos, l\'imites y overflow}

\begin{itemize}
  \item \texttt{int}: aproximadamente $[-2\cdot10^9,2\cdot10^9]$.
  \item \texttt{long long}: aproximadamente $[-9\cdot10^{18},9\cdot10^{18}]$.
  \item \texttt{double}: unas 15 cifras decimales de precisi\'on.
  \item \texttt{long double}: precisi\'on dependiente de la plataforma.
  \item \texttt{\_\_int128}: entero no est\'andar disponible en GCC/Clang.
\end{itemize}

Promover \emph{antes} de operar:

\begin{lstlisting}[style=cpp]
ll promoted_product(int a, int b) { return 1LL * a * b; }
bool square_leq(ll x, ll n) { return x <= n / x; } // x > 0

void print_i128(i128 x) {
    if (x == 0) { cout << 0; return; }
    bool negative = x < 0;
    u128 value = negative ? (u128)(-(x + 1)) + 1 : (u128)x;
    if (negative) cout << '-';
    string s;
    while (value) {
        s.push_back(char('0' + value % 10));
        value /= 10;
    }
    reverse(s.begin(), s.end());
    cout << s;
}
\end{lstlisting}

El overflow de enteros con signo es comportamiento indefinido en C++.
